% probe02: 风险名字候选（电位器/极性电容/齐纳/LED/battery2）——nonstopmode 批量试错
\documentclass[margin=5pt]{standalone}
\usepackage[american]{circuitikz}
\begin{document}
\begin{circuitikz}
  \draw (0,8) to[pR] ++(2,0);       \node[anchor=west,font=\ttfamily] at (2.3,8) {pR};
  \draw (6,8) to[cC] ++(2,0);       \node[anchor=west,font=\ttfamily] at (8.3,8) {cC};
  \draw (0,6) to[eC] ++(2,0);       \node[anchor=west,font=\ttfamily] at (2.3,6) {eC};
  \draw (6,6) to[zD] ++(2,0);       \node[anchor=west,font=\ttfamily] at (8.3,6) {zD};
  \draw (0,4) to[leD] ++(2,0);      \node[anchor=west,font=\ttfamily] at (2.3,4) {leD};
  \draw (6,4) to[battery2] ++(2,0); \node[anchor=west,font=\ttfamily] at (8.3,4) {battery2};
  % 电位器 wiper anchor 试用
  \draw (0,2) to[pR, name=P1] ++(2,0); \draw[red] (P1.wiper) -- ++(0,0.8);
  \node[anchor=west,font=\ttfamily] at (2.3,2) {pR.wiper};
\end{circuitikz}
\end{document}
